\chapter*{คำนำ}
\addcontentsline{toc}{chapter}{คำนำ}

\textbf{คู่มือปฏิบัติการฟิสิกส์สำหรับจุลชีววิทยา} (\textbf{Physics for Microbiology Laboratory Manual}) เล่มนี้ จัดทำขึ้นเพื่อใช้ประกอบการเรียนการสอนภาคปฏิบัติการของรายวิชา \textbf{4011111 ฟิสิกส์สำหรับจุลชีววิทยา} จำนวน 3(2-3-4) หน่วยกิต (บรรยาย 2 ชั่วโมง และปฏิบัติการ 3 ชั่วโมงต่อสัปดาห์) สำหรับนักศึกษาปริญญาตรี หลักสูตรวิทยาศาสตรบัณฑิต สาขาวิชาจุลชีววิทยา คณะวิทยาศาสตร์และเทคโนโลยี มหาวิทยาลัยราชภัฏรำไพพรรณี โดยใช้ควบคู่กับเอกสารคำสอนเล่มหลักในชุดเดียวกัน

แนวคิดหลักในการออกแบบคู่มือปฏิบัติการเล่มนี้ คือ การเปลี่ยนภาพจำของปฏิบัติการฟิสิกส์แบบดั้งเดิมที่เน้นสมการคณิตศาสตร์ซับซ้อน ให้กลายเป็น \textbf{ปฏิบัติการฟิสิกส์ที่รับใช้ห้องปฏิบัติการจุลชีววิทยาโดยตรง} ทั้ง 13 การทดลองในเล่มนี้ได้รับการออกแบบให้เชื่อมโยงกับเครื่องมือวิทยาศาสตร์และปรากฏการณ์ทางชีวภาพที่นักจุลชีววิทยาต้องพบเจอในวิชาชีพจริง อาทิ การสอบเทียบความแม่นยำของไมโครปิเปตด้วยเครื่องชั่งดิจิทัล การสังเกตการเคลื่อนที่แบบบราวน์เทียบกับการว่ายเคลื่อนที่ของแบคทีเรีย การวัดความหนืดและอัตราเร็วตกตะกอนตามกฎของสโตกส์ การคำนวณค่าแรงเหวี่ยงสัมพัทธ์ ($\text{RCF}$ หรือ $\times g$) และการจัดสมดุลโรเตอร์ของเครื่องปั่นเหวี่ยง (\textbf{Centrifuge}) การตรวจสอบเสถียรภาพอุณหภูมิของตู้อบเลี้ยงเชื้อ (\textbf{Incubator}) การวิเคราะห์ความดันไอและจุดเดือดในหม้อนึ่งความดันไอ (\textbf{Autoclave}) ที่สภาวะ $121\,^\circ\text{C}$ ($15\text{ psi}$) กลไกประจุไฟฟ้าในการย้อมสีแกรม การแยกโมเลกุลด้วยสนามไฟฟ้าในเจลอะกาโรส (\textbf{Gel Electrophoresis}) การใช้เครื่องวัดดัชนีหักเห (\textbf{Refractometer}) ติดตามการหมักน้ำตาลของยีสต์ บทบาทของน้ำมันช่วยส่อง (\textbf{Immersion Oil}) ในกล้องจุลทรรศน์กำลังขยาย $1{,}000\times$ ประสิทธิภาพการฆ่าเชื้อของรังสีอัลตราไวโอเลต (\textbf{UV-C}) ไปจนถึงการวัดความหนาแน่นเซลล์ที่ความยาวคลื่น $600\text{ nm}$ ($\text{OD}_{600}$) ด้วยเครื่องสเปกโตรโฟโตมิเตอร์ (\textbf{Spectrophotometer})

เพื่อให้สอดคล้องกับพื้นฐานคณิตศาสตร์ของนักศึกษา คู่มือปฏิบัติการแต่ละบทจึงประกอบด้วยสรุปสูตรสำเร็จรูป ตัวอย่างการคำนวณข้อมูลล่วงหน้าทีละขั้นตอน ขั้นตอนการทดลองที่ชัดเจน และ\textbf{ใบบันทึกผลการทดลองประจำกลุ่ม}ที่นักศึกษาสามารถกรอกข้อมูลจริง วาดกราฟ และตอบคำถามท้ายการทดลองได้ทันที ผู้เขียนหวังเป็นอย่างยิ่งว่าคู่มือปฏิบัติการเล่มนี้จะช่วยสร้างทักษะปฏิบัติการที่แม่นยำ ปลอดภัย และมีจรรยาบรรณทางวิทยาศาสตร์ให้แก่นักศึกษาสาขาจุลชีววิทยาทุกคน

\vspace{1.0cm}
\begin{flushright}
    \textbf{ผู้ช่วยศาสตราจารย์ ดร.ชีวะ ทัศนา}\\[0.2cm]
    สาขาวิชาฟิสิกส์ คณะวิทยาศาสตร์และเทคโนโลยี\\
    มหาวิทยาลัยราชภัฏรำไพพรรณี\\
    ภาคเรียนที่ 2 ปีการศึกษา 2569
\end{flushright}
